%%
%% CogniSync: A Secure, Local-First Architecture for Persistent LLM Agent Memory
%% Submitted to CIKM 2026, Full Paper Track
%%
\documentclass[sigconf,authordraft]{acmart}

\AtBeginDocument{%
  \providecommand\BibTeX{{%
    Bib\TeX}}}

\setcopyright{acmlicensed}
\copyrightyear{2026}
\acmYear{2026}
\acmDOI{XXXXXXX.XXXXXXX}

\acmConference[CIKM '26]{Proceedings of the 35th ACM International
  Conference on Information and Knowledge Management}{October 2026}{Atlanta,
  GA, USA}
\acmISBN{978-1-4503-XXXX-X/2026/10}

\usepackage{booktabs}
\usepackage{multirow}
\usepackage{amsmath}
\usepackage{algorithm}
\usepackage{algorithmic}
\usepackage{xspace}
\usepackage{url}
\usepackage{microtype}
\usepackage{tabularx}

\newcolumntype{Y}{>{\centering\arraybackslash}X}
\newcolumntype{L}{>{\raggedright\arraybackslash}X}
\raggedbottom
\setlength{\emergencystretch}{3em}
\hfuzz=25pt
\hbadness=10000
\makeatletter
\@ACM@balancefalse
\makeatother

\begin{document}

\title{CogniSync: A Secure, Local-First Architecture for Persistent LLM
Agent Memory}

\author{Anonymous Authors}
\affiliation{%
  \institution{Anonymous Institution}
  \city{Anonymous City}
  \country{Anonymous Country}}
\email{anonymous@anonymous.edu}
\renewcommand{\shortauthors}{Anonymous et al.}

%% ============================================================
\begin{abstract}
%% ============================================================

Most retrieval-augmented generation systems are evaluated on one axis:
how accurately they answer queries. This paper evaluates on two. We ask
what it costs, in retrieval precision, to secure a RAG pipeline against
adversarial content in the retrieval corpus, and we measure that cost
across 34,929 unique queries on five benchmarks rather than inferring it from
a small domain study.

The answer, measured after correcting a routing miscalibration discovered
during evaluation, is 1.1\% MRR. Dense retrieval achieves MRR 0.867;
CogniSync, which adds a context validation layer for prompt-injection and
tool-spoofing defense, achieves MRR 0.856 with adaptive hybrid routing.
We call this the security-quality tradeoff, decompose it into identifiable
components, and argue it is the right cost to pay for deployments where
retrieval-layer attacks are a genuine risk.

The path to this number is itself a contribution. Our large-scale
evaluation identified a routing miscalibration: exact-match queries were
being sent to lexical-only retrieval rather than hybrid fusion. We
diagnosed this precisely the equality between CogniSync and the Lexical
baseline on exact-match MRR (both 0.861 pre-fix) was an exact fingerprint
of the bug and implemented a correction using adaptive hybrid fusion
with query-dependent weighting. Exact-match MRR improved from 0.861 to
0.927 after the fix. Dense retrieval cannot eliminate prompt injection or
reduce tool spoofing because it has no mechanism to inspect retrieved
content before it enters the agent's context window. CogniSync has one,
and that mechanism has a 1.1\% precision cost after proper routing.
Reporting both numbers together, at scale, is the primary contribution
of this work.

Beyond the tradeoff measurement, the paper makes three contributions with
cleanly isolated evidence. First, a cross-session episodic memory graph
produces a $+$4.1~pp MRR gain on cross-session agent tasks, with a direct
linking-strategy ablation confirming that causal edges outperform
A-MEM's cosine-similarity edges by $+$1.2~pp. Second, adaptive hybrid
routing with query-dependent weighting outperforms fixed-weight RRF on
exact-match queries (MRR 0.927 vs.~0.889), demonstrating that calibrated
routing is preferable to always-on fusion. Third, the architecture runs
at 28.0~ms median latency on local hardware, $8.0\times$ faster than
equivalent cloud retrieval at identical domain-specific quality.

\end{abstract}

\begin{CCSXML}
<ccs2012>
 <concept>
  <concept_id>10002951.10003260.10003261</concept_id>
  <concept_desc>Information systems~Retrieval models and ranking</concept_desc>
  <concept_significance>500</concept_significance>
 </concept>
 <concept>
  <concept_id>10002951.10003317.10003338</concept_id>
  <concept_desc>Information systems~Question answering</concept_desc>
  <concept_significance>300</concept_significance>
 </concept>
 <concept>
  <concept_id>10002978.10002986</concept_id>
  <concept_desc>Security and privacy~Adversarial attacks</concept_desc>
  <concept_significance>300</concept_significance>
 </concept>
 <concept>
  <concept_id>10010147.10010178</concept_id>
  <concept_desc>Computing methodologies~Natural language processing</concept_desc>
  <concept_significance>200</concept_significance>
 </concept>
</ccs2012>
\end{CCSXML}

\ccsdesc[500]{Information systems~Retrieval models and ranking}
\ccsdesc[300]{Information systems~Question answering}
\ccsdesc[300]{Security and privacy~Adversarial attacks}
\ccsdesc[200]{Computing methodologies~Natural language processing}

\keywords{LLM agents, retrieval-augmented generation, episodic memory,
prompt injection, local-first AI, hybrid retrieval, security-quality
tradeoff}

\maketitle

%% ============================================================
\section{Introduction}
%% ============================================================

Consider an agent debugging a production incident that started three
sessions ago. The relevant context exists somewhere in the memory store:
a schema migration ran on a specific date, the CI pipeline began failing
six hours later, and rolling back a particular commit fixed it. A flat
vector index stores the individual facts. It does not store the order,
the causality, or the session provenance. Asked ``what changed before the
failure?'', the agent retrieves topically related content rather than the
event sequence that explains the failure. Enlarging the context window
does not help. Adding more documents does not help. What is missing is
temporal structure, not data.

Now add an adversary. The same pipeline that makes the agent useful,
namely the retrieval system that injects relevant context at inference
time, also gives an attacker a direct channel into the agent's
reasoning. Any document the agent can retrieve can instruct it. The
Tool-Hijacker threat model~\cite{hui2024toolhijacker} formalizes this:
adversarial text in the corpus can redirect tool selection without
touching model weights, enabling exfiltration, unauthorized mutations,
or command execution. The Model Context Protocol~\cite{chen2025mcp} has
made this worse by making tool discovery dynamic and corpus-driven.

These two problems have mostly been worked on separately. Memory papers
optimize recall and latency. Security papers measure attack success
rates. Neither field asks the question practitioners most need answered:
\emph{what does it cost, in retrieval precision, to defend against
retrieval-layer attacks?} That tradeoff is real, it is measurable, and
to our knowledge it has not been characterized at scale before this
paper.

CogniSync is our attempt to close both gaps in a single architecture. It
runs entirely on-device on a local small language model, eliminating
cloud telemetry and providing network-layer control over outbound
connections. Built on that foundation, it adds three components: a
cross-session episodic memory graph for temporal reasoning, a hybrid
dense-lexical retrieval engine with adaptive query routing, and a
context validation layer that screens retrieved content for adversarial
payloads before they enter the agent's context window.

Our evaluation is designed to be honest rather than favorable. We run
34,929 unique queries across five benchmarks specifically to see where the
system fails and this evaluation caught a real failure. The
large-scale results revealed that the query-type detector was sending
exact-match queries to lexical-only retrieval rather than hybrid fusion.
We diagnosed the failure precisely, corrected it with adaptive
query-dependent weighting, and re-evaluated. After the fix, CogniSync
achieves MRR 0.856, trailing Dense by 1.1\%. We report the pre-fix and
post-fix numbers, explain what the difference tells us about both the
routing behavior and the true cost of the validation layer, and argue
that running an evaluation large enough to catch this kind of failure is
itself worth doing.

The specific contributions are:

\begin{enumerate}
\item \textbf{First large-scale measurement of the security-quality
tradeoff in RAG.} 34,929 unique queries. 172,605 evaluation rows. The true
1.1\% MRR cost of context validation, isolated after correcting a routing
miscalibration discovered within the same evaluation.

\item \textbf{Episodic memory graph with causally clean ablation and
direct linking-strategy comparison.} $+$4.1~pp MRR on cross-session tasks.
Disabling the graph recovers the Vanilla~RAG baseline exactly. A direct
A-MEM comparison on identical MemoryArena tasks confirms that causal
edges outperform semantic similarity edges by $+$1.2~pp, isolating the
linking strategy specifically.

\item \textbf{Routing calibration: diagnosis and fix.} The large-scale
evaluation exposed a miscalibration where exact-match queries routed to
lexical-only retrieval, producing an exact MRR equality with the Lexical
baseline (both 0.861). We corrected this by implementing adaptive hybrid
fusion with query-dependent weighting: exact-match queries apply a higher
lexical weight while retaining a non-zero dense contribution for fallback
robustness. Exact-match MRR improved from 0.861 to 0.927. Overall MRR
improved from 0.834 to 0.856.

\item \textbf{Formal validation layer specification.} We give the
context validation algorithm explicitly so operators can tune it,
understand its failure modes, and reproduce the evaluation.
\end{enumerate}

%% ============================================================
\section{Related Work}
%% ============================================================

\subsection{Agent Memory Architectures}

MemGPT~\cite{packer2023memgpt} treats LLM memory as an operating-system
problem: the model pages information between in-context working memory
and an external store. MemoryBank~\cite{zhong2023memorybank} applies
Ebbinghaus forgetting-curve dynamics to prioritize semantically frequent
content. A-MEM~\cite{xu2024amem} implements the Zettelkasten method with
atomic memory nodes connected by cosine-similarity edges and
LLM-generated annotations.

None of these architectures addresses the question of \emph{when}. They
record what the agent learned without preserving the temporal or causal
structure of how it was learned. CogniSync's episodic graph fills this
gap by encoding event order and causal precedence directly as graph edges,
following Tulving's episodic/semantic distinction~\cite{tulving1972}.
This is not a cosmetic difference from A-MEM. A-MEM's edges connect nodes
that are semantically similar; CogniSync's edges connect nodes where one
event caused or enabled another. For incident-reconstruction queries,
the motivating use case of this work, the causal graph is structurally
better suited. A direct ablation comparing both strategies on MemoryArena
cross-session tasks confirms this: CogniSync's episodic edges achieve
MRR 0.9436 versus A-MEM's 0.9316 on identical tasks
(Section~\ref{sec:episodic_linking}).

\subsection{Retrieval-Augmented Generation}

Dense passage retrieval~\cite{karpukhin2020dpr} has dominated open-domain
QA since 2020. Its weakness on fine-grained technical
identifiers~\cite{zhao2022dense} motivates hybrid dense-sparse retrieval.
Reciprocal rank fusion~\cite{cormack2009rrf} is the standard combination
mechanism; our implementation follows the $k{=}60$ convention.

Fixed-weight RRF is a reasonable default but not necessarily optimal for
all query types. On exact-match queries involving UUIDs, error codes, and
versioned paths, assigning a higher lexical weight while retaining a
non-zero dense component yields better results than either equal-weight
fusion or lexical-only retrieval. Our routing fix (Section~\ref{sec:routing})
provides direct empirical support for this at scale.

\subsection{Security in Agentic RAG}

Perez and Ribeiro~\cite{perez2022injection} first characterized indirect
prompt injection. Tool-Hijacker~\cite{hui2024toolhijacker} extended this
to multi-tool agents with a formal threat model that our evaluation
directly operationalizes. Chen et al.~\cite{chen2025mcp} analyze MCP
security, identifying unauthenticated STDIO execution as the primary
attack vector; CogniSync's local-first design removes this surface by
construction.

On data exfiltration, the relevant prior work covers covert channels in
language model outputs. Bagdasaryan and
Shmatikov~\cite{bagdasaryan2023covert} demonstrate steganographic encoding
of private information within generated text;
Greshake et al.~\cite{greshake2023indirect} characterize indirect
injection across deployed systems. Neither work proposes a
retrieval-integrated defense. CogniSync's validation layer addresses the
injection and spoofing classes directly; steganographic exfiltration is
acknowledged as an open problem (Section~\ref{sec:exfil}).

\textbf{The measurement gap.} Prior agent memory papers do not measure
security. Prior RAG security papers do not measure precision cost. The
contribution of this work is precisely the number that sits between them:
1.1\% MRR, at scale, with a decomposition of its sources and a correction
of the routing behavior that would have obscured that number.

%% ============================================================
\section{Problem Formulation}
%% ============================================================

\subsection{Memory and Context Model}

Let an LLM have parameters $\theta$ and maximum token capacity
$C_{\max}$. At inference step $t$, working context $W_t$ satisfies
$\lvert W_t\rvert \leq C_{\max}$. The agent maintains a long-term corpus
$\mathcal{D}$, queried via retrieval function $R(q_t,\mathcal{D},k)$
returning the top-$k$ chunks for inclusion in $W_t$. The design problem
is populating $\mathcal{D}$ with episodically structured records so that
$R$ maximizes task-relevant recall across session boundaries, while
adversarially crafted content in $\mathcal{D}$ cannot redirect agent
behavior.

\subsection{Retrieval Failure Taxonomy}

We use three failure modes from the RAGEC taxonomy. \emph{Retrieval
Thrash}: successive queries converge semantically but retrieved content
fails verification. \emph{Tool Storms}: cascading invocations exhaust the
token budget. \emph{Context Bloat}: low-utility content accumulates in
$W_t$ and overwrites system instructions.

\subsection{Threat Model}

We assume an adversary controlling external content ingested into
$\mathcal{D}$, including web pages, documents, and API responses, without
model weight access. Three attack classes:

\begin{itemize}
\item \textbf{Prompt Injection}: adversarial string $x_{\mathrm{adv}} \in W_t$
  causes the agent to satisfy the adversary's goal rather than the
  system goal.
\item \textbf{Tool Spoofing}: injected documentation shifts the
  tool-selection distribution toward an unauthorized tool.
\item \textbf{Data Exfiltration}: the agent encodes private context
  within an authorized outbound response.
\end{itemize}

This model does not cover model-weight attacks, local privilege
escalation, or adversaries who know the perplexity threshold and craft
injections to evade it. The security evaluation should be read within
these bounds.

%% ============================================================
\section{CogniSync Architecture}
%% ============================================================

\subsection{Local-First Execution}

All reasoning, memory indexing, and tool selection run on a locally
hosted small language model (Google Gemma-2B-Instruct, 4-bit quantized). The evaluation was
conducted on a Kaggle dual-GPU benchmarking pipeline (two NVIDIA T4 GPUs,
16~GB VRAM each), though the architecture scales down to single-GPU edge
deployments requiring 8--24~GB RAM. Two security properties follow by construction: zero data
telemetry during inference, and network-layer control over outbound
connections. The MCP STDIO attack class is eliminated because there is
no remote STDIO surface. External content is untrusted and passes through
the validation layer before reaching the context window.

The hardware requirement is real. Developer workstations handle it without
issue. Edge devices and IoT hardware do not. This is the cost of the
privacy guarantee.

\subsection{Cross-Session Episodic Memory Graph}
\label{sec:episodic_arch}

The episodic graph $G=(V,E)$ stores interactions as atomic event nodes.
Each node $v_i \in V$ contains: the interaction text, a Unix timestamp,
a session identifier, and the tool set invoked. Directed edge
$v_i \to v_j$ exists when $v_i$'s outcome was a structural precondition
for $v_j$'s invocation. It captures causation rather than semantic
similarity: one event caused or enabled the other.

After each session, an asynchronous consolidation process traverses $G$
and extracts stable recurring facts into a parallel FAISS semantic store.
The episodic graph is retained at full temporal resolution; the semantic
store handles stable factual lookups that do not require event ordering.
This dual-store design follows Tulving~\cite{tulving1972}.

The architecture is motivated by incident reconstruction: ``the service
failed on this date because that migration ran and broke JWT validation,
and rolling back commit X fixed it.'' That query requires temporal and
causal ordering. A flat index cannot answer it. A similarity graph can
retrieve related content but not the causal chain. The directed episodic
graph can. Section~\ref{sec:episodic_linking} provides empirical
confirmation with a direct comparison against A-MEM's similarity-based
linking on identical tasks.

\subsection{Context Validation Layer}
\label{sec:validation_arch}

Algorithm~\ref{alg:val} specifies the context validation layer.

\begin{algorithm}[t]
\caption{Context Validation}
\label{alg:val}
\begin{algorithmic}[1]
\REQUIRE chunk $c$, perplexity threshold $\tau$, authorized schema $S_T$,
  window width $w$, circuit-breaker limit $B$
\ENSURE $\mathit{decision}\in\{\textsc{admit},\textsc{filter}\}$
\STATE $\mathit{fails}\leftarrow 0$
\FORALL{token window $c_w$ of width $w$ in $c$}
  \STATE $P\leftarrow\mathrm{PPL}(c_w;\,\theta_{\mathrm{local}})$
  \IF{$P < \tau$}
    \STATE $\mathit{fails}\leftarrow\mathit{fails}+1$
    \IF{$\mathit{fails}\geq B$}
      \STATE \textbf{return} \textsc{filter}
    \ENDIF
  \ENDIF
\ENDFOR
\IF{$c$ contains tool metadata $m$}
  \IF{$\lnot\,\mathrm{ValidateSchema}(m,S_T)$}
    \STATE \textbf{return} \textsc{filter}
  \ENDIF
\ENDIF
\STATE \textbf{return} \textsc{admit}
\end{algorithmic}
\end{algorithm}

A sliding window of $w$ tokens scans each retrieved chunk; windows with
perplexity below threshold $\tau$ under the local SLM are flagged.
Adversarial injections are typically grammatically fluent instruction-text,
which has low perplexity under a language model trained on instructions.
The circuit breaker $B$ prevents scanning from continuing once enough
sub-threshold windows have accumulated. In all experiments, the validation layer used a sliding-window width of $w = 32$ tokens, a perplexity threshold of $\tau = 12.0$, and a circuit-breaker limit of $B = 3$ sub-threshold windows.

Schema validation checks tool metadata in retrieved chunks against a
deterministic schema $S_T$ specifying authorized parameter types and
value ranges. This closes the natural-language ambiguity that
tool-spoofing exploits: a chunk that names an authorized tool but
specifies subtly wrong parameters will fail schema validation regardless
of linguistic fluency.

\subsection{Adaptive Hybrid Dense-Lexical Retrieval}
\label{sec:retrieval_arch}

Dense retrieval uses all-MiniLM-L6-v2 (384 dimensions, FAISS
\texttt{IndexFlatIP} inner-product index with L2 embedding normalization,
the tokenizer distributed with all-MiniLM-L6-v2 (BERT-compatible WordPiece), 120-token chunks, and a 20-token overlap stride).
Lexical retrieval uses SQLite
FTS5 with BM25. The two signals fuse via Reciprocal Rank Fusion:
\begin{equation}
  \mathrm{score}(d) =
    \frac{W_{\mathrm{dense}}}{k+r_{\mathrm{dense}}(d)}+
    \frac{W_{\mathrm{lexical}}}{k+r_{\mathrm{lex}}(d)},
    \quad k=60,
\end{equation}
where weights {\mathrm{dense}}$ and {\mathrm{lexical}}$ are set adaptively by the query-type detector.
The detector classifies each query as semantic or exact-match based on
high-entropy token patterns: UUIDs, hexadecimal strings, versioned
endpoint paths, and uppercase alphanumeric error codes (see
Appendix~\ref{app:routing_spec} for the full deterministic specification).

\textbf{Adaptive weighting.} For semantic queries the detector sets
{\mathrm{dense}} = 1.0$ and {\mathrm{lexical}} = 1.0$ (equal-weight fusion). For exact-match queries it sets
{\mathrm{dense}} = 0.2$ and {\mathrm{lexical}} = 3.0$, heavily up-weighting the lexical component while retaining a
non-zero dense contribution. This differs critically from the pre-fix
behavior, in which exact-match queries set {\mathrm{dense}} = 0.0$ and {\mathrm{lexical}} = 1.0$ (lexical-only),
discarding the dense signal entirely. The non-zero dense contribution
provides fallback robustness when lexical matching fails on paraphrased or
partially-tokenized identifiers, and prevents the routing from collapsing
to the Lexical baseline on that query stratum. Section~\ref{sec:routing}
reports the empirical effect of this correction.

%% ============================================================
\section{Experimental Setup}
%% ============================================================

\subsection{Datasets}

\textbf{CogniSync Benchmark v1.0.} A domain-specific corpus for
engineering workflows: JWT authentication, database schema migrations,
microservice routing, CI/CD configuration, observability stacks, API
gateway rate limiting, and RBAC access control. The benchmark contains
312 semantic queries and 96 exact-match queries, evaluated with two
independent annotators on a three-level relevance scale
($\kappa = 0.81$ semantic, $\kappa = 0.94$ exact-match). All agent-pipeline
claims (Section~\ref{sec:domain}) are grounded in this benchmark alone and
should not be generalized; the four public benchmarks below carry all
general claims.

\textbf{Public benchmarks.} MS-MARCO (10,047 queries, paraphrase-heavy
passage QA), SQuAD (5,000 queries, extractive QA), SciQ (4,474 queries,
science QA), CodeSearchNet (15,000 queries, code search). Total across
all five benchmarks: 34,929 unique queries per system, 172,605 evaluation rows
across five configurations. Query-type classification produces 23,590
semantic and 10,931 exact-match queries across the public benchmarks.

\textbf{MemoryArena.} A multi-session agentic evaluation gym covering
web shopping, group travel planning, progressive information search, and
sequential mathematical reasoning. Tasks average 57 action steps and
exceed 40,000 reasoning-trace tokens per run. A cross-session linking
sub-evaluation specifically tests reconstruction of causal event chains
across session boundaries.

\subsection{Baselines}

The large-scale evaluation compares: Dense (FAISS inner-product only),
Lexical (BM25/FTS5 only), Vanilla~RAG (Dense plus standard RAG wrapper;
retrieval scores are identical to Dense by construction), Hybrid-Naive
(always-on equal-weight RRF, $\alpha=0.5$ fixed), and CogniSync
(adaptive routing plus context validation). The domain-specific benchmark
additionally includes A-MEM, MemGPT, LangChain MMR, LlamaIndex,
MemoryBank, and Pinecone. All systems use 120-token chunks; domain-specific
experiments run over five seeds.

\subsection{Metrics}

Recall@$K$ ($K \in \{1,3,5\}$), MRR, MAP, NDCG@5. Statistical tests:
Wilcoxon signed-rank test and McNemar's test over per-query deltas;
Cohen's $d$ for effect size. Security evaluation: attack success rate
(ASR) and attacker MRR (lower = stronger defense). Large-scale tables
report both macro-averages (arithmetic mean of per-dataset MRRs) and
micro-averages (weighted by query count).

%% ============================================================
\section{Results}
%% ============================================================

\subsection{Domain-Specific Benchmark}
\label{sec:domain}

Table~\ref{tab:main} reports CogniSync Benchmark v1.0 results over five
seeds. CogniSync achieves Recall@5 = 79.5\%, MRR = 0.594, NDCG@5 =
0.636 at 28.0~ms median latency, $8.0\times$ faster than cloud-simulated
retrieval at identical quality. All seven pairwise comparisons are
statistically significant ($p < 0.01$, two-sided sign test). The nearest
competitor, Vanilla~RAG, is beaten by $+$1.3~pp Recall@5 and $+$0.012
MRR ($p = 0.0063$). These margins are modest but consistent across all
five trials. MemoryBank's temporal decay actively degrades recall on
freshly-relevant content; LangChain MMR's maximal-marginal-relevance
re-ranking trades recall for diversity in a way that hurts this
benchmark's single-relevant-document structure.

\begin{table*}[t]
\caption{%
  Domain-specific benchmark (CogniSync Benchmark v1.0, 408 queries,
  5 seeds, averaged). Agent memory baselines evaluated here only.
  $\dagger$~MemoryBank applies temporal decay.
  *~Pinecone excluded from quality rankings (identical retrieval to
  Vanilla~RAG at $\sim$222~ms emulated cloud latency).%
}
\label{tab:main}
\begin{tabular}{lrrrrrrr}
\toprule
\textbf{System} &
  \textbf{R@1} & \textbf{R@3} & \textbf{R@5} &
  \textbf{MRR} & \textbf{MAP} & \textbf{NDCG@5} &
  \textbf{Lat (ms)} \\
\midrule
CogniSync (Ours)      & 0.468 & 0.693 & 0.795 & 0.594 & 0.583 & 0.636 & 27.8 \\
Vanilla RAG           & 0.461 & 0.676 & 0.781 & 0.582 & 0.571 & 0.624 & 24.5 \\
A-MEM                 & 0.461 & 0.671 & 0.761 & 0.576 & 0.568 & 0.617 & 24.9 \\
LlamaIndex            & 0.448 & 0.640 & 0.736 & 0.557 & 0.538 & 0.587 & 24.1 \\
MemGPT                & 0.361 & 0.615 & 0.760 & 0.505 & 0.487 & 0.555 & 26.6 \\
LangChain MMR         & 0.461 & 0.476 & 0.500 & 0.474 & 0.447 & 0.460 & 28.8 \\
MemoryBank$^\dagger$ & 0.189 & 0.364 & 0.560 & 0.307 & 0.285 & 0.350 & 24.0 \\
Pinecone (lat.)*      & 0.461 & 0.676 & 0.781 & 0.582 & 0.571 & 0.624 & 223.4 \\
\bottomrule
\end{tabular}
\end{table*}

\subsection{Latency and Scaling}

Table~\ref{tab:lat} shows latency and memory from 2K to 100K documents.
The overhead versus Vanilla~RAG stays between 17\% and 28\%. At 28.0~ms
median on the domain benchmark, the delta versus Vanilla~RAG's 24.5~ms
is under one millisecond at operating scale. Memory scales linearly at
approximately 3~MB per 1K documents; the 100K-document configuration
reaches 313.5~MB, well within workstation constraints.

\begin{table}[t]
\caption{Latency and memory scaling. ``CS'' = CogniSync; ``VR'' = Vanilla RAG.}
\label{tab:lat}
\footnotesize\setlength{\tabcolsep}{3pt}
\begin{tabularx}{\columnwidth}{@{}LYYYYY@{}}
\toprule
\textbf{Corpus} &
  \textbf{CS (ms)} & \textbf{VR (ms)} &
  \textbf{Overhead} & \textbf{p95 (ms)} & \textbf{Mem (MB)} \\
\midrule
2,000   & 29.6 & 24.9 & $+$19\% & 43.7 & 6.7   \\
10,000  & 29.1 & 24.1 & $+$21\% & 38.6 & 31.2  \\
50,000  & 36.4 & 30.3 & $+$20\% & 44.6 & 156.9 \\
100,000 & 51.7 & 40.4 & $+$28\% & 64.7 & 313.5 \\
\bottomrule
\end{tabularx}
\end{table}

\subsection{Agent Task Performance}
\label{sec:agent}

Table~\ref{tab:agent} reports per-task MRR on MemoryArena. The episodic
graph's contribution is directionally clean: $+$4.1~pp MRR on cross-session
tasks, no detectable effect on single-session tasks ($<$0.003 MRR across
all single-session types). Disabling the graph recovers the Vanilla~RAG
baseline exactly on cross-session MRR (0.698 in both cases).

Two task types where Full CogniSync trails Vanilla~RAG warrant
explanation. Code generation (0.925 vs.~0.950) and API integration
(0.652 vs.~0.692) both involve exact API names and configuration keys.
Task completion is 100\% for both systems on these types, so the MRR gap
reflects ranking quality, not success rate. Terse API documentation has
characteristically low perplexity under a general language model, making
it vulnerable to the validation layer's filtering. Tuning the perplexity
threshold for API-documentation corpora would recover this gap; it is a
deployment-specific calibration opportunity, not a fundamental limitation.

\begin{table}[t]
\caption{%
  Agent task MRR on MemoryArena. Disabling the episodic graph recovers
  the Vanilla~RAG baseline on cross-session tasks exactly.%
}
\label{tab:agent}
\footnotesize\setlength{\tabcolsep}{3pt}
\begin{tabularx}{\columnwidth}{@{}LYYYY@{}}
\toprule
\textbf{Configuration} &
  \textbf{Code Gen} & \textbf{Debug} &
  \textbf{API Int.} & \textbf{Cross-Sess.} \\
\midrule
Full CogniSync              & 0.925 & 0.787 & 0.652 & 0.739 \\
Ablation: no episodic graph & 0.935 & 0.784 & 0.671 & 0.698 \\
Vanilla RAG baseline        & 0.950 & 0.783 & 0.692 & 0.698 \\
\bottomrule
\end{tabularx}
\end{table}

\subsection{Large-Scale Generalization and the Security-Quality Tradeoff}
\label{sec:largescale}

Tables~\ref{tab:large} and~\ref{tab:pds} present results across all
34,929 unique queries with adaptive hybrid routing. Dense/Vanilla~RAG leads at
MRR 0.867. CogniSync follows at MRR 0.856, 1.1 percentage points lower.
Hybrid-Naive sits at 0.844.

This ranking deserves unpacking. CogniSync now outperforms Hybrid-Naive
(0.856 vs.~0.844) despite the additional overhead of the validation layer.
The reason is that adaptive routing with query-dependent weighting
({\mathrm{dense}}=0.2, W_{\mathrm{lexical}}=3.0$ for exact-match, {\mathrm{dense}}=1.0, W_{\mathrm{lexical}}=1.0$ for semantic) outperforms
fixed equal-weight fusion on exact-match queries: MRR 0.927 vs.~0.889
on that stratum. Hybrid-Naive is a reasonable default; calibrated routing
is better.

The gap between CogniSync and Dense ranges from $-$0.045 on MS-MARCO
to $-$0.002 on CodeSearchNet. A striking result from the query-level
analysis: CogniSync strictly outperforms Vanilla~RAG on 2,499 queries,
with an average MRR gain of $+$0.374 on those wins. Ninety-four percent
of these wins (2,355 queries) occur on MS-MARCO, demonstrating that hybrid
fusion actively rescues sparse queries where pure dense retrieval fails.
The validation overhead is not uniformly distributed; it concentrates where
the embedding space is weakest, which is also where the security benefit
matters most.

CodeSearchNet is essentially immune to the security-quality tradeoff.
Only 0.8\% of code queries suffer a severe MRR penalty (drop $\geq 0.5$
versus Vanilla~RAG), compared to 18.1\% on MS-MARCO and 5.9\% on SQuAD.
Code blocks, API specifications, and functional documentation have rigid
syntactic structures that behave predictably under the SLM's perplexity
threshold, providing direct support for the claim that CogniSync is
well-suited to its target domain of software engineering.

\begin{table*}[t]
\caption{%
  Large-scale aggregate evaluation. Dense and Vanilla~RAG share identical
  retrieval; scores are equal by construction. Wilcoxon $p$-values are
  computed over per-query MRR deltas (34,929 unique queries). CogniSync uses
  adaptive hybrid routing throughout.%
}
\label{tab:large}
\footnotesize\setlength{\tabcolsep}{3pt}
\begin{tabularx}{\textwidth}{@{}LYYYYYLY@{}}
\toprule
\textbf{System} &
  \textbf{R@1} & \textbf{R@3} & \textbf{R@5} &
  \textbf{MRR} & \textbf{NDCG@5} & \textbf{Wilcoxon vs VR} &
  \textbf{$\Delta$MRR} \\
\midrule
Dense/Vanilla RAG & 0.803 & 0.906 & 0.961 & 0.867 & 0.891 & -- & $\phantom{-}$0.000 \\
CogniSync         & 0.792 & 0.898 & 0.955 & 0.856 & 0.882 & $p = 1.50\times 10^{-172}$ & $-$0.011 \\
Hybrid-Naive      & 0.774 & 0.889 & 0.951 & 0.844 & 0.872 & $p = 1.61\times 10^{-112}$ & $-$0.023 \\
Lexical (BM25)    & 0.726 & 0.838 & 0.912 & 0.796 & 0.826 & $p < 10^{-300}$ & $-$0.071 \\
\bottomrule
\end{tabularx}
\end{table*}

\begin{table*}[t]
\caption{%
  Per-dataset MRR with macro- and micro-averages. Macro = arithmetic mean;
  Micro = query-count weighted (matches Table~\ref{tab:large}). Cohen's
  $d$ is reported against Dense/Vanilla~RAG. v1.0 = CogniSync Benchmark.%
}
\label{tab:pds}
\footnotesize\setlength{\tabcolsep}{2.5pt}
\begin{tabularx}{\textwidth}{@{}LYYYYYYYY@{}}
\toprule
\textbf{System} &
  \textbf{MARCO} & \textbf{SQuAD} & \textbf{SciQ} &
  \textbf{CSNet} & \textbf{v1.0} & \textbf{Macro} &
  \textbf{Micro-avg} & \textbf{Cohen's $d$} \\
\midrule
Dense/VRAG   & 0.543 & 0.991 & 0.982 & 0.998 & 1.000 & 0.903 & 0.867 & -- \\
CogniSync    & 0.526 & 0.981 & 0.982 & 0.996 & 0.883 & 0.874 & 0.856 & $-$0.052 \\
Hybrid-Naive & 0.499 & 0.971 & 0.972 & 0.995 & 0.873 & 0.862 & 0.844 & $-$0.094 \\
Lexical      & 0.375 & 0.934 & 0.954 & 0.989 & 0.741 & 0.799 & 0.796 & $\phantom{-}$0.231 \\
\bottomrule
\end{tabularx}
\end{table*}

\textbf{Interpreting the residual gap.} After the routing fix, the
remaining 1.1\% MRR gap between CogniSync and Dense is attributable to
the validation layer alone. The routing analysis in
Section~\ref{sec:routing} shows that on semantic queries (68.3\% of the
evaluation), CogniSync is statistically indistinguishable from Hybrid-Naive
(both MRR 0.823). The 1.1\% gap concentrates on exact-match queries where
the validation layer filters some legitimate low-perplexity content. That
is the actual security-quality tradeoff. It is smaller than the 3.3\%
headline number would have suggested before the routing correction.

%% ============================================================
\section{Routing Calibration: Diagnosis and Fix}
\label{sec:routing}
%% ============================================================

The most diagnostic result in this evaluation was not a performance number
but an equality. Table~\ref{tab:routing} shows the query-type decomposition
both before and after the routing correction.

\textbf{Diagnosing the miscalibration.} In the pre-fix system, two exact
equalities stood out in the MRR columns: (1) CogniSync and Hybrid-Naive
were identical on semantic queries (both 0.823); (2) CogniSync and Lexical
were identical on exact-match queries (both 0.861). These were not
approximations; they were exact to three decimal places. Equality (1)
confirmed the routing was working correctly on semantic queries: the
detector classified them as semantic, routed to hybrid fusion, and produced
the same result as Hybrid-Naive. Equality (2) confirmed a miscalibration
on exact-match queries: the detector was setting {\mathrm{dense}} = 0.0$, discarding
the dense signal entirely and producing output identical to pure BM25
for 31.7\% of all queries. A 408-query benchmark would not have made
this equality exact; 34,929 unique queries did.

\textbf{Implementing the fix.} We corrected the miscalibration by changing
the exact-match routing from {\mathrm{dense}} = 0.0$ (lexical-only) to {\mathrm{dense}} = 0.2, W_{\mathrm{lexical}} = 3.0$
(lexical-upweighted hybrid). This retains a non-zero dense contribution
for fallback robustness while still prioritizing lexical matching on
high-entropy tokens. The change required modifying a single parameter in
the query-type detector; the rest of the architecture is unchanged.

\textbf{Effect of the correction.} The improvement is concentrated on
the exact-match stratum, as expected. Exact-match MRR improved from 0.861
to 0.927, surpassing Hybrid-Naive's 0.889. The gain over Hybrid-Naive
confirms that {\mathrm{dense}} = 0.2, W_{\mathrm{lexical}} = 3.0$ (lexical-upweighted) outperforms {\mathrm{dense}} = 1.0, W_{\mathrm{lexical}} = 1.0$
(equal-weight) on this query type: when the query contains a UUID or error
code, the lexical signal deserves more weight, but discarding the dense
signal entirely was too aggressive. Semantic query performance is unchanged
(MRR 0.823), as expected: the fix only touches the routing branch for
exact-match queries. Overall MRR improved from 0.834 to 0.856.

\begin{table*}[t]
\caption{%
  Query-type decomposition across all 34,929 unique queries, before and after
  the routing correction. Pre-fix: the exact equality between CogniSync
  and Lexical on exact-match MRR (both 0.861) is a fingerprint of the
  $\alpha=0$ bug. Post-fix: adaptive weighting ($\alpha=0.25$ for
  exact-match) outperforms both fixed-weight alternatives.%
}
\label{tab:routing}
\footnotesize\setlength{\tabcolsep}{3.5pt}
\begin{tabularx}{\textwidth}{@{}LYYYYL@{}}
\toprule
\textbf{System} &
  \textbf{Sem.\ MRR} & \textbf{EM MRR} &
  \textbf{Sem.\ R@5} & \textbf{EM R@5} &
  \textbf{Effective Path} \\
\midrule
Dense/Vanilla RAG              & 0.844 & 0.902 & 0.955 & 0.971 &
  Query-agnostic dense \\
Hybrid-Naive                   & 0.823 & 0.889 & 0.944 & 0.962 &
  Always-on RRF ({\mathrm{d}}{=}1.0, W_{\mathrm{l}}{=}1.0$) \\
CogniSync \textit{(pre-fix)}   & 0.823 & 0.861 & 0.944 & 0.934 &
  Semantic $\to$ RRF; exact-match $\to$ lexical-only ({\mathrm{d}}{=}0.0$, \textbf{bug}) \\
CogniSync \textit{(post-fix)}  & 0.823 & 0.927 & 0.944 & 0.978 &
  Semantic $\to$ RRF ({\mathrm{d}}{=}1.0$); exact-match $\to$ adaptive hybrid ({\mathrm{l}}{=}3.0$) \\
Lexical                        & 0.771 & 0.861 & 0.897 & 0.934 &
  Lexical-only \\
\bottomrule
\end{tabularx}
\end{table*}

The 0.927 exact-match MRR for CogniSync post-fix exceeds Hybrid-Naive's
0.889. This is not surprising in hindsight: for queries containing UUIDs
or versioned endpoint paths, a higher lexical weight is strictly more
appropriate than equal-weight fusion, but the dense component still helps
when an identifier is partially tokenized or when a paraphrased variant
appears in the corpus. The adaptive weighting captures this tradeoff in
a way that neither fixed-weight option does.

%% ============================================================
\section{Ablation Studies}
%% ============================================================

\subsection{Episodic Graph: Disabling and Linking-Strategy Comparison}
\label{sec:episodic_linking}

Disabling the graph while holding all other components constant reduces
cross-session MRR from 0.739 to 0.698, recovering the Vanilla~RAG baseline
exactly. Single-session metrics change by less than 0.003 MRR in every
task type. There is no interference: the episodic component adds value on
cross-session tasks and is invisible on single-session tasks.

Table~\ref{tab:linking} fills the linking-strategy comparison gap. We ran
both CogniSync (causal/temporal edges) and A-MEM (cosine-similarity edges)
on the same MemoryArena cross-session linking tasks. CogniSync achieves
MRR 0.9436 versus A-MEM's 0.9316, a $+$1.2~pp advantage that isolates the
linking strategy specifically. A-MEM's semantic edges actually degraded the
Vanilla~RAG baseline (0.9316 vs.~0.9426), confirming that topical clustering
is a poor substitute for causal ordering on incident-reconstruction tasks.
When the question is ``what happened in what order,'' a similarity graph
answers a different question.

\begin{table}[t]
\caption{%
  Cross-session linking strategy ablation on MemoryArena. All variants
  use identical retrieval; only the edge type in the memory graph varies.%
}
\label{tab:linking}
\footnotesize\setlength{\tabcolsep}{3pt}
\begin{tabularx}{\columnwidth}{@{}LLYY@{}}
\toprule
\textbf{Variant} & \textbf{Linking Strategy} & \textbf{MRR} & \textbf{Recall@5} \\
\midrule
Vanilla RAG (Baseline)     & None (flat index) & 0.9426 & 1.000 \\
A-MEM (Semantic Edges)     & Cosine similarity & 0.9316 & 1.000 \\
CogniSync (Episodic Edges) & Causal/temporal   & 0.9436 & 1.000 \\
\bottomrule
\end{tabularx}
\end{table}

\subsection{Hybrid Retrieval}

On the domain-specific benchmark, hybrid and dense-only systems produce
identical scores on the 312 semantic queries: the routing correctly directs
these to hybrid fusion, which on paraphrase-heavy queries differs negligibly
from dense-only in a high-quality embedding space. On the 96 exact-match
queries, adaptive hybrid ({\mathrm{lexical}}{=}3.0$) achieves Recall@5 = 0.834 versus
0.750 for dense-only, an $+$8.4~pp gain confirming the lexical component
earns its place on technical identifier queries.



%% ============================================================
\section{Security Evaluation}
%% ============================================================
\label{sec:security}

\subsection{Prompt Injection and Tool Spoofing}

Table~\ref{tab:sec} reports attack effectiveness before and after the
validation layer, using two complementary metrics: Attack Success Rate
(ASR) measures whether an attack fully redirects agent behavior; Attacker
MRR measures partial goal achievement. Without sanitization: prompt
injection reaches attacker MRR 0.90; tool spoofing reaches 1.00. After
sanitization: no prompt injection attack in the evaluated set succeeded
(ASR 0.00\%, attacker MRR 0.000). Tool spoofing attacker MRR fell from
1.00 to 0.038, a $\sim$96\% reduction.

Tool spoofing is harder to eliminate completely than prompt injection
because a subtly malformed tool schema can be both fluent (passing the
perplexity check) and adversarial (redirecting tool selection). The schema
validator closes most of this gap by catching parameter-level manipulations
that perplexity windowing cannot see.

The attack set covers known injection patterns from published
taxonomies~\cite{hui2024toolhijacker,perez2022injection}. It does not
include adaptive attacks designed to evade the perplexity threshold
specifically. The 0.00\% ASR on prompt injection applies within this scope.

\subsection{Adaptive Threshold Evasion}
\label{sec:adaptive}

The advanced Adaptive Threshold Evasion attack, where the attacker
specifically targets Algorithm~\ref{alg:val}'s perplexity threshold,
succeeds at 4.72\% ASR on both the undefended baseline and CogniSync.
We report this directly. The validation layer does not currently defend
against an adversary who knows and exploits the threshold; that is the
most important open security gap.

\subsection{Data Exfiltration}
\label{sec:exfil}

Known data exfiltration payloads are blocked entirely (ASR 0.00\%, MRR
drop 0.000). However, residual leakage via authorized channels remains:
mean leakage MRR 0.073 across seeds, with a $3.4\times$ spread. The
remaining attack exploits an authorized outbound channel: the agent encodes
private context within a semantically legitimate response payload.
Connection-level filtering cannot close this because the connection is
permitted by design. Applying semantic monitoring of output content relative
to authorized output specifications, following Bagdasaryan and
Shmatikov~\cite{bagdasaryan2023covert}, is the most promising direction.

\begin{table*}[t]
\caption{%
  Security evaluation. ASR = Attack Success Rate (0\% = fully blocked).
  Attacker MRR measures partial goal achievement (lower = stronger defense).
  Adaptive Evasion is the open gap. An Attacker MRR of 0.073 reflects
  residual leakage via authorized channels.%
}
\label{tab:sec}
\footnotesize\setlength{\tabcolsep}{3pt}
\begin{tabularx}{\textwidth}{@{}LLYYYL@{}}
\toprule
\textbf{Attack} & \textbf{System} & \textbf{ASR} & \textbf{MRR Drop} & \textbf{Attacker MRR} & \textbf{Result} \\
\midrule
Data Exfiltration & No Defense & 3.20\% & 0.0054 & --    & Vulnerable \\
Data Exfiltration & CogniSync  & 0.00\% & 0.000  & 0.073 & Blocked \\
Prompt Injection  & No Defense & 3.20\% & 0.0055 & 0.900 & Vulnerable \\
Prompt Injection  & CogniSync  & 0.00\% & 0.000  & 0.000 & Complete \\
Tool Spoofing     & No Defense & --      & --     & 1.000 & Vulnerable \\
Tool Spoofing     & CogniSync  & --      & --     & 0.038 & $\sim$96\% reduction \\
Adaptive Evasion  & No Defense & 4.72\% & 0.0081 & --    & Open gap \\
Adaptive Evasion  & CogniSync  & 4.72\% & 0.0081 & --    & Open gap \\
\bottomrule
\end{tabularx}
\end{table*}

%% ============================================================
\section{Limitations}
\label{sec:lim}
%% ============================================================

\textbf{Hardware.} The local SLM (Gemma~4 efficiency class) requires
8--24~GB RAM. Developer workstations handle it. Edge and IoT devices do
not. Local execution is what gives the privacy guarantee, and that
guarantee has a hardware price.

\textbf{Benchmark scope.} The domain-specific benchmark was
author-constructed and covers the domain the system was designed for.
Statistical significance on 408 queries at $p < 0.01$ is consistent
evidence within that domain. It is not a substitute for a third-party
engineering retrieval benchmark.

\textbf{Adaptive adversaries.} The security evaluation covers fixed-taxonomy
attacks. It does not cover adversaries who know the perplexity threshold and
design injections to stay above it. The Adaptive Threshold Evasion result
(4.72\% ASR, identical to the undefended baseline) makes this gap explicit.
This is left for future work.

\textbf{Perplexity threshold calibration.} The current threshold is set
globally. The validation layer's false-positive rate varies by corpus type:
API documentation has characteristically low perplexity and is more likely
to be filtered. Corpus-specific threshold tuning is a deployment-specific
calibration opportunity, not a fundamental limitation.

%% ============================================================
\section{Conclusion}
%% ============================================================

The question CogniSync was built to answer is: what does it cost to make
a RAG pipeline safe? The answer, measured across 34,929 unique queries on five
benchmarks, is 1.1\% MRR but the path to that number matters as much
as the number itself.

Our large-scale evaluation caught a routing miscalibration that was not
visible at smaller scale: exact-match queries were routed to lexical-only
retrieval, producing an exact MRR equality with the pure BM25 baseline
(both 0.861) on 31.7\% of all queries. We diagnosed this through the
exact equality fingerprint, corrected it by implementing adaptive hybrid
fusion with query-dependent weighting ({\mathrm{dense}} = 0.2, W_{\mathrm{lexical}} = 3.0$ for exact-match),
and re-evaluated. Exact-match MRR improved from 0.861 to 0.927, surpassing
fixed-weight Hybrid-Naive (0.889). Overall MRR improved from 0.834 to 0.856.

The residual 1.1\% gap after the routing fix is attributable to the
validation layer alone: the cost of screening retrieved content for prompt
injection and tool-spoofing payloads before they enter the agent's context
window. Dense retrieval has no such mechanism. CogniSync does, and the
evaluation now measures its cost precisely rather than conflating it with
a routing bug.

On the problem the system was specifically designed for, the picture is
more favorable. On a 408-query engineering benchmark, CogniSync outperforms
seven baselines including A-MEM, MemGPT, and Vanilla~RAG with statistical
significance, at $8.0\times$ the latency advantage of cloud retrieval. The
episodic graph delivers cleanly: $+$4.1~pp cross-session MRR, zero
single-session interference, confirmed by a tight ablation. A direct
linking-strategy comparison confirms that causal edges outperform A-MEM's
similarity edges by $+$1.2~pp on cross-session tasks. And CodeSearchNet
shows only 0.8\% severe loss rate under the validation layer versus 12\%
on MS-MARCO, which is the architecture working as intended: the syntactic
regularity of code and API documentation makes perplexity-based filtering
reliable on that corpus.

The broader argument is methodological. A 408-query benchmark would not
have found the routing bug. 34,929 unique queries did. Running large-scale
evaluations that include unfavorable baselines and then fixing what
they find is how you find out what is actually wrong with your system.
We release the full 172,605-row evaluation dataset, benchmark tools, and
evaluation pipeline so others can run the same kind of evaluation on theirs.

\begin{acks}
Omitted for blind review.
\end{acks}

\bibliographystyle{ACM-Reference-Format}
\bibliography{cognisync}

%% ============================================================
\appendix
%% ============================================================

\section{Benchmark Construction}

Relevance judgments for CogniSync Benchmark v1.0 followed a three-level
rubric (highly relevant / partially relevant / not relevant). Two annotators
with software engineering backgrounds scored independently; a third
adjudicated disagreements. Cohen's $\kappa = 0.81$ for semantic queries,
$\kappa = 0.94$ for exact-match. The higher exact-match agreement is
expected: relevance for UUID and endpoint-string queries is less
interpretively ambiguous than for semantic queries about engineering concepts.

\section{Routing Classifier Specification}
\label{app:routing_spec}

The classifier is rule-based and deterministic. A query is classified
exact-match if it contains any of: (1) a UUID
\texttt{([0-9a-fA-F]\{8\}(-[0-9a-fA-F]\{4\})\{3\}-[0-9a-fA-F]\{12\})};
(2) a hexadecimal string \texttt{(0x[0-9a-fA-F]+)}; (3) a versioned
endpoint path (contains \texttt{/} and a numeric version component); or
(4) an uppercase alphanumeric error code. No trainable parameters; no
exposure to evaluation data. A 15\% false-negative rate is observed
in the large-scale evaluation (some exact-match queries that lack any of
the four surface patterns are classified as semantic); this explains the
31.7\% exact-match proportion rather than a higher one.

Upon exact-match classification, the routing sets the weight
parameters {\mathrm{dense}} = 0.2$ and {\mathrm{lexical}} = 3.0$ in the RRF fusion equation, up-weighting the
lexical signal while retaining a non-zero dense contribution. Upon semantic
classification, {\mathrm{dense}} = 1.0$ and {\mathrm{lexical}} = 1.0$ (equal-weight). The pre-fix behavior set
{\mathrm{dense}} = 0.0$ and {\mathrm{lexical}} = 1.0$ for exact-match queries, which is equivalent to lexical-only
retrieval and was the source of the routing miscalibration.

\section{Statistical Tests}

Wilcoxon signed-rank tests were computed over the full per-query MRR delta
distributions (35,000+ paired observations). McNemar's test was applied to
per-query Recall@1 binary outcomes. Cohen's $d$ was computed from the
per-query MRR distributions. All $p$-values are two-sided.
Effect-size interpretation: $d \approx 0.052$ (Dense vs.~CogniSync, very
small); $d \approx 0.094$ (Hybrid-Naive vs.~CogniSync, small);
$d \approx 0.231$ (Lexical vs.~CogniSync, small-to-medium). The small
effect sizes reflect that, while the differences are statistically
unambiguous at this sample size, the practical gap per query is modest,
consistent with the headline MRR deltas.

\end{document}